% Generated draft: requires mathematical and source review.
\documentclass[11pt]{article}
\usepackage[margin=25mm]{geometry}
\usepackage{fontspec}
\setmainfont{Latin Modern Roman}
\usepackage{amsmath,amssymb,mathtools,unicode-math}
\setmathfont{Latin Modern Math}
\usepackage{enumitem,needspace,xurl,hyperref}
\setlength{\parindent}{0pt}
\setlength{\parskip}{5pt}
\setlength{\emergencystretch}{3em}
\newcommand{\N}{\mathbb{N}}
\newcommand{\Z}{\mathbb{Z}}
\newcommand{\Q}{\mathbb{Q}}
\newcommand{\R}{\mathbb{R}}
\newcommand{\C}{\mathbb{C}}
\newcommand{\F}{\mathbb{F}}
\newcommand{\A}{\mathbb{A}}
\newcommand{\PP}{\mathbb{P}}
\newcommand{\E}{\mathbb{E}}
\newcommand{\eps}{\varepsilon}
\newcommand{\dd}{\,\mathrm{d}}
\begin{document}
\section*{UnsolvedMath ID: 700814}
\textbf{Source problem code:} catalogue entry
\subsection{Definitions and mathematical statement}
A group \(G\) is torsion-free if \(g^n=1\), for \(g\in G\) and an integer \(n\geq 1\), implies \(g=1\). For a specified coefficient ring \(A\) with identity, let \(R=A[G]\) denote its group ring: the ring of finite formal sums \(\sum_{g\in G}a_g g\), with multiplication determined by \((ag)(bh)=(ab)(gh)\).

Using left modules as a convention, a free \(R\)-module is a module isomorphic to \(\bigoplus_{i\in I}R\) for some set \(I\), with no finiteness restriction on \(I\). A free cyclic submodule of a module \(M\) is a submodule \(Rm\), generated by an element \(m\in M\), for which the map \(R\to Rm\), \(r\mapsto rm\), is an isomorphism. Equivalently, \(\{r\in R:rm=0\}=\{0\}\).

For the coefficient ring intended by the source, the question is whether the following assertion holds: for every torsion-free group \(G\), every set \(I\), and every nonzero left \(A[G]\)-submodule \(M\leq\bigoplus_{i\in I}A[G]\), there exists \(m\in M\setminus\{0\}\) such that
\[
\forall r\in A[G],\qquad rm=0\ \Longrightarrow\ r=0.
\]
The supplied statement does not specify the coefficient ring or quantify over coefficient rings. Thus \(A\) in this formulation is an unresolved parameter, not an additional universal quantifier. If the intended group ring is the integral group ring, the displayed assertion is to be read with \(A=\Z\); this interpretation requires confirmation. The source also leaves module sidedness unstated. For right modules, the corresponding formulation uses \(mR\) and the condition \(mr=0\Longrightarrow r=0\).
\subsection{Short English statement}
The question asks whether every nonzero submodule of a free module over the group ring of a torsion-free group contains a submodule isomorphic to that group ring as a module over itself. The coefficient ring and module sidedness are unspecified in the supplied source.
\subsection{Sources}
\begin{itemize}
\item Local UnsolvedMath registry, supplied complete source record: permanent ID \(700814\), exact source problem code \texttt{catalogue entry}, title \emph{Kourovka Notebook Problem 11.122}. The registry labels the problem \emph{open}; this is supplied metadata, not an independently verified assessment.
\item Attributed to A. L. Shmel'kin, proposed in 1990. The supplied background cites \emph{The Kourovka Notebook (Unsolved Problems in Group Theory)}, 11th issue (1990), Problem 11.122, and the 21st issue edition (2026), \texttt{arXiv:1401.0300v48}. The supplied collection description names E. I. Khukhro and V. D. Mazurov as editors.
\item The supplied registry record was created and updated at \texttt{2026-10-16T14:00:00Z}, with upstream revision \texttt{d371ee917ac9647fbb06cd6fc529deec6f3dda82}. These are source-record timestamps, not independent literature-verification dates. No external source was consulted for this draft.
\end{itemize}
\end{document}
